% AGE — manuscript for Nature Machine Intelligence (Article format)
% Compile: pdflatex manuscript.tex  (figures in ./fig)

\documentclass[11pt,a4paper]{article}

\usepackage[margin=2.4cm]{geometry}
\usepackage{graphicx}
\usepackage{amsmath,amssymb}
\usepackage{booktabs}
\usepackage[super,sort&compress]{natbib}
\usepackage[colorlinks=true,urlcolor=blue,citecolor=blue]{hyperref}
\usepackage[expansion=false]{microtype}
\usepackage{enumitem}
\usepackage{xcolor}

\setlength{\parskip}{0.4em}
\linespread{1.12}

\newcommand{\Age}{\textsc{Age}}
\newcommand{\PhysFormer}{\textsc{PhysFormer}}
\newcommand{\pn}{\textsuperscript}

\title{\Large\bfseries An artificial general engineer that learns software and
physical engineering through verified trial and error}

\author{
  Sehaj Randhir Singh$^{1*}$ \\
  \small\textit{$^{1}$Independent Researcher; affiliated with the Department of
  Electrical and Computer Engineering, NYU Tandon School of Engineering}\\
  \small$^{*}$Correspondence: sehajrsinghs@gmail.com
}

\date{}

\begin{document}

\maketitle
\vspace{-1.2em}

\begin{abstract}
\noindent Engineering is one of the last frontiers of artificial intelligence:
it demands not prediction but closed-loop action, verification against reality,
and memory across attempts. Commercial efforts (e.g.\ Prometheus, the \$12B
venture founded by Jeff Bezos and Vik Bajaj) are pursuing this goal as an
``artificial general engineer'' for physical products, but no complete,
minimal, and verifiable instantiation of the concept has been made public. We
present \Age, an open autonomous engineering agent that learns through
trial and error gated by verification: a planner (deterministic or
language-model driven) selects actions from a sandboxed skill layer; every
critical step is verified against tests or \emph{independent} numeric
simulations before it is accepted; failures are journaled as lessons and
replayed into future plans. The physics core introduces \PhysFormer, a
transformer whose reasoning layers are trained jointly with
physics-consistency layers that penalize violations of the governing
equations (a physics-informed, PINN-style residual), and, as its central
methodological contribution, a new way to feed physics into a transformer:
\emph{Law-Conditioned Attention (LCA)}. Each governing equation is
tokenized into a fixed 22-operator symbolic vocabulary
(\texttt{first\_time\_deriv}, \texttt{second\_space\_deriv},
\texttt{self\_advection}, \texttt{moment\_curvature},
\texttt{energy\_conservation}, \ldots), embedded into a law vector, and
injected into every attention layer and into the shared readout, so
attention is conditioned on the physics in force rather than merely
penalized by it after the fact. Inputs are labeled by physical quantity
from a shared vocabulary, so the simply supported beam and the cantilever
present \emph{identical} parameter tokens and only the law signature
disambiguates them; one shared transformer then solves six physical laws
(beam, cantilever, projectile, pendulum, spring--mass, RC) at once, with a
paired six-seed test against a constant-signature ablation showing the
equation signal significantly improves trajectory fidelity
($p=0.0003$ over 36 runs, 21\% median curve-error reduction) and reduces
held-out answer error on exactly the laws whose parameters do not identify
them (beam 34\% vs.\ 40\%, cantilever 28\% vs.\ 35\%), while an honest
negative result (RC, whose normalized trajectory is the same curve for
every time constant) localizes where conditioning does not pay. We further
show the equation signal is causal, not merely predictive: swapping the
law signature at inference on \emph{identical} inputs displaces the model's
prediction across the boundary between the two bending solutions
(steering index $+0.090$ vs.\ $-0.132$ for the constant-signature
ablation, $p<0.0001$; per-seed $p=0.031$) while the control is provably
insensitive,and the same generalist adapts to a law held out of pretraining with 25\%
of the specialist's data ($2.4$--$2.9\times$ lower answer error than a
from-scratch specialist on the same samples). The two physics channels are
then cleanly separated: the residual-loss channel cuts the predicted
field's governing-equation violation 4--8$\times$ at equal data without
improving field accuracy, while the input (LCA) channel is what buys
trajectory fidelity and adaptation. Across all
eight physics domains (projectile, pendulum, spring--mass, simply
supported beam, cantilever beam, viscous Burgers conservation law, RC
circuit, and steady-state heat conduction on a square plate), closed-form
answers are verified to relative residuals below $5\times10^{-2}$ (often
$10^{-9}$) by independent integrators. On identical budgets (256 samples, 60
epochs, five seeds, per-seed weight initialization), a pooled Wilcoxon
signed-rank test over 25 paired runs shows the physics term reduces the
predicted trajectories' violation of the governing equations
significantly ($p<0.001$, Cliff's $\delta=-0.92$, a $\sim$19-fold median
reduction) relative to an equal-capacity no-physics transformer, while
leaving scalar accuracy statistically unchanged ($p=0.65$)---including on
the nonlinear Burgers equation, whose shock profile the physics-
consistency layer keeps physical---and it matches or beats both that
baseline and an equal-capacity MLP on scalar accuracy wherever the
parameter space is rich enough for the extra machinery to pay off (the
simplest two-parameter domain is an honest exception). A verification-gate
benchmark across seven missions shows the gate eliminates false successes
(0\% versus 29\% without it). The complete system
passes 54 unit tests and completes end-to-end engineering missions
(scaffold a project and run its tests; design a beam, cantilever, or
Burgers flow and verify it against independent numerics) in a single
planning iteration. \Age is an existence proof that the ``general
engineer'' loop---plan, act, verify, reflect, remember---can be built today,
in roughly 5,900 lines of dependency-light code, and it identifies the
specific open problems (trajectory-level physics consistency, verifier-
driven repair) that must be solved to scale the idea to real hardware.
\end{abstract}

\noindent\textbf{Keywords:} autonomous agents; physics-informed machine learning;
trial-and-error learning; engineering automation; transformers

\section{Introduction}

Modern machine learning excels at \emph{prediction}: given a static input,
produce a static output---an image label, a sentence, a scalar
property\pn{\citep{lecun2015,brown2020,silver2017}}. Engineering, by
contrast, is an \emph{interactive} discipline. An engineer does not simply
predict whether a beam will fail; she designs it, simulates it, builds it,
measures it, and revises it when the measurement disagrees with the model.
The defining loop of engineering is therefore closed-loop trial and error in
which the \emph{world}---a test suite, a simulator, or a physical
instrument---is the final arbiter of correctness. Sutton's ``bitter lesson''
argues that systems which exploit computation and feedback outperform those
that encode brittle human knowledge\pn{\citep{sutton2019}}; engineering is
precisely the regime in which that feedback is expensive, sparse, and
physically grounded.

Large language models (LLMs) have recently been embedded in agent loops that
interleave reasoning and action\pn{\citep{yao2023,wei2022}}, that reflect on
their own failures\pn{\citep{shinn2023}}, that explore open-ended
environments\pn{\citep{wang2023,park2023}}, and that repair real software
repositories\pn{\citep{jimenez2024,li2022}}. A parallel line of work in
physics-informed machine learning has shown that embedding governing
equations as soft constraints dramatically improves learning and
generalization on forward and inverse problems in
mechanics\pn{\citep{raissi2019,karniadakis2021,lu2021,chen2018}}. Yet these
two lines have largely proceeded separately: agents verify weakly (if at
all), and physics-informed networks predict but do not act.

In June 2026, Bezos and Bajaj publicly described Prometheus, a company
founded on the explicit thesis of an ``artificial general engineer''
(\Age): artificial intelligence that learns engineering through real-world
trial and error, aimed at physical products such as aerospace, automotive,
and computing hardware, funded at \$12B with a \$41B
valuation\pn{\citep{techcrunch2026,wikipedia2025}}. The ambition is clear;
what has been missing is a public, complete, minimal instantiation of the
underlying architecture---small enough to inspect, honest enough to show its
failure modes, and general enough to span both software and physics.

Here we present \Age, such an instantiation. \Age is a complete autonomous
engineering agent built from four components: (i) a \emph{planner} that
turns a goal into a sequence of skill invocations (a deterministic
``mechanical'' brain with zero dependencies, or a general LLM brain for any
OpenAI-compatible endpoint); (ii) an \emph{embodiment layer} of sandboxed
skills (inspect, read, search, edit, run, verify) confined to a working
directory; (iii) a \emph{verification gate}: any step marked
\texttt{expect: ok} must pass tests or an independent numeric simulation, or
the mission fails and a lesson is journaled; and (iv) an \emph{episodic
journal} that records every episode and replays lessons into future plans.
For the physical world, \Age ships with \PhysFormer, a transformer whose
reasoning layers are trained \emph{together} with physics-consistency layers
that penalize violations of the governing equations, following the
physics-informed paradigm\pn{\citep{raissi2019}}. Every physics answer is
additionally cross-checked by a numerically independent simulator (Euler or
RK4 integration, finite-difference beam solver), so that the agent's
verification gate is never a tautology of its own formulas.

We report three classes of results. First, the agent loop is robust: 54 unit
tests cover mission completion, the verification gate catching injected
bugs, sandbox confinement, and journal memory, and a gate benchmark shows
the gate eliminates false successes (0\% vs 29\% for a verifier-free
baseline on the same missions). Second, the physics core is accurate and
self-consistent: across eight domains, closed-form answers agree with
independent numerics to relative residuals typically below $10^{-3}$
(the canonical beam design agrees to $2.6\times10^{-9}$); \PhysFormer with
the physics term beats equal-budget MLP and no-physics baselines on scalar
accuracy where the parameter space is large enough for attention to matter,
and on trajectory-level consistency with the governing equations everywhere
(one to two orders of magnitude lower residual, including on the nonlinear
Burgers equation); the physics residual drops by one to two orders of
magnitude during training. Third, an end-to-end mission---``design a simply
supported beam with span 4 m, point load 3 kN, modulus $2\times10^{11}$ Pa,
inertia $5\times10^{-6}$~m$^4$''---is parsed from natural language, solved
by closed form, verified by finite differences to $2.6\times10^{-9}$, and
predicted by \PhysFormer to within a few percent of the exact answer in a
single planning iteration; identical missions for a cantilever and a
Burgers flow complete through the same single-step plan.

\section{Results}

\subsection{Architecture: the verification-gated trial-and-error loop}

Figure~\ref{fig:arch} shows the system. A mission begins with a goal in
natural language (``scaffold a Python calculator project and make its tests
pass'', ``design a simply supported beam carrying 3\,kN over 4\,m''). The
planner---either the deterministic mechanical brain or an LLM---produces an
ordered plan of skill invocations. Each step is executed by the embodiment
layer inside a sandboxed working directory, and steps annotated
\texttt{expect: ok} act as gates: if a verification step fails, the mission
stops, a lesson is derived from the failure output, and both the episode and
the lesson are appended to the journal (a JSONL file). Subsequent missions
receive journal lessons as context, so the agent's behavior improves across
episodes without retraining---an episodic analogue of the memory--consolidation
trade-off studied in continual learning\pn{\citep{kirkpatrick2017}}.

\begin{figure}[t]
\centering
\includegraphics[width=0.98\textwidth]{fig/fig1_architecture.png}
\caption{\textbf{System architecture.} The planner (mechanical or LLM brain)
emits skill invocations executed by the sandboxed embodiment layer. The
verification gate (\texttt{expect: ok}) requires tests or independent numeric
simulations to pass before a mission succeeds; failures write lessons to the
episodic journal, which is replayed into future plans.}
\label{fig:arch}
\end{figure}

The loop is deliberately minimal (Algorithm~\ref{alg:loop}): plan, act,
verify, reflect, remember. Its power comes from where the authority lies.
Unlike agents that treat the model's own judgment as the arbiter of
success\pn{\citep{yao2023,shinn2023}}, \Age delegates correctness to
\emph{execution}: a test runner, a typechecker, or an independent physics
simulator. This is the same design choice that made verifier-based training
effective for mathematical reasoning\pn{\citep{cobbe2021,lightman2024}} and
competitive code generation\pn{\citep{li2022}}, here applied to the full
engineering loop rather than to a single generation step.

\begin{figure}[t]
\centering
\includegraphics[width=0.95\textwidth]{fig/fig2_physformer.png}
\caption{\textbf{\PhysFormer: a transformer adjusted for physics.}
Problem parameters are embedded as (parameter-id, value) token pairs and
processed by reasoning layers (a standard transformer encoder). Two heads
decode a scalar engineering answer and a reference trajectory. A
physics-consistency layer then evaluates how far the predicted trajectory
strays from the governing equation, and that residual enters the loss, so
the network is trained on data and physics \emph{together}.}
\label{fig:physformer}
\end{figure}

\subsection{The verification gate catches injected faults}

A central design claim is that the verification gate is load-bearing, not
decorative. We tested it by injection: after scaffolding a pristine Python
project whose four tests pass, we silently corrupted the implementation
(\texttt{return a + b} $\rightarrow$ \texttt{return a - b}). The agent
planned the mission ``run the test suite and check it passes''; the verify
step failed; the mission was correctly reported as \emph{failed}; and a
lesson (\texttt{verify step failed: ...}) was written to the journal. The
same gate refuses to scaffold over a non-empty directory (protecting
existing work) and blocks destructive shell commands
(\texttt{rm -rf /}, \texttt{git push}, \texttt{sudo}) by default, with a
path-traversal guard on every file operation. These properties are enforced
by dedicated tests rather than asserted informally.

\subsection{Two brains: deterministic and learned planning}

The mechanical brain recognizes mission types from goal text (scaffold,
explain, scan for TODO markers, run the test suite, solve a physics design)
and emits deterministic plans; it requires no API key, no network, and no
pretrained weights, which makes the entire loop testable in unit tests. The
LLM brain is a drop-in replacement that re-plans each iteration from
observations, enabling genuinely novel plans (e.g.\ editing code to fix a
failing test) at the cost of an external endpoint. Both brains share the
same embodiment, verification gate, and journal; the choice is a single
flag. This two-brain design decouples the \emph{loop} (which we argue is the
scientifically interesting object) from the \emph{planner} (which is
improving rapidly in the wider LLM ecosystem).

\subsection{The physics core: eight domains, closed forms, and independent verification}

The physics layer (\texttt{physx}, $\sim$3,300 lines of PyTorch) implements
eight domains---projectile motion, pendulum, spring--mass oscillator, simply
supported beam under a central point load, cantilever beam under a tip
load, the viscous Burgers conservation law, RC circuit charging, and
steady-state heat conduction on a square plate---each
with three artifacts: an exact closed-form solution (the training target), a
reference trajectory sampled from that solution (the data), and an
\emph{independent} numeric simulator used for verification. Crucially, the
numeric verifiers are written without reference to the closed forms: Euler
integration for the projectile and RC circuit, RK4 for the pendulum,
finite-difference solution of the fourth-order beam equation (the
cantilever verifier integrates the moment distribution $M(x)=-P(L-x)$
from static equilibrium twice with a ghost point at the clamped end),
finite-volume upwind integration of the Burgers conservation form, a
sparse finite-difference Poisson solve for the heat plate, and an
ODE
residual for the spring. Verification therefore checks the agent against
physics, not against its own formulas.

Table~\ref{tab:verify} reports the agreement between closed form and numeric
simulation across domains. All residuals are at least two orders of
magnitude below the tolerance threshold for ``verified'' (the beam example
used throughout this paper agrees to $2.6\times10^{-9}$).

\begin{table}[t]
\centering
\small
\begin{tabular}{lccl}
\toprule
Domain & Answer (canonical example) & Numeric method & Relative residual \\
\midrule
Projectile & range $=15.90$ m & Euler (2,000 steps) & $<2\times10^{-3}$ \\
Pendulum & period $=2.09$ s & RK4 ($\Delta t=10^{-4}$) & $<10^{-3}$ \\
Spring--mass & $\omega=\sqrt{k/m}$ & ODE residual, central diff. & $<10^{-3}$ \\
Beam & max deflection $=4.0$ mm & Finite difference (400 nodes) & $2.6\times10^{-9}$ \\
Cantilever & max deflection $=64$ mm & Finite difference (400 nodes) & $3.1\times10^{-6}$ \\
Burgers & peak $u=1.23$ m/s & Finite volume, upwind (2,000 cells) & $<5\times10^{-2}$ \\
RC circuit & $\tau=RC=1.0$ s & Euler (4,000 steps) & $<2\times10^{-3}$ \\
Heat2D & peak $T=300$ K & Sparse FD Poisson (48$^2$ grid) & $<2\times10^{-3}$ \\
\bottomrule
\end{tabular}
\caption{\textbf{Closed-form answers verified by numerically independent
simulators.} Tolerances for the ``verified'' flag are $2\times10^{-3}$ for
projectile and RC, $10^{-3}$ for pendulum and spring, $2\times10^{-2}$
for the beam and cantilever finite-difference solvers, and $5\times10^{-2}$
for the Burgers finite-volume solver (shock cases push the coarse-grid
residual up; at 2{,}000 cells the two solvers agree to $\sim$1\%).}
\label{tab:verify}
\end{table}

\subsection{\PhysFormer: reasoning layers and physics layers, together}

\PhysFormer (Figure~\ref{fig:physformer}) maps a problem's parameters to
both a scalar engineering answer and a trajectory. Parameters are
normalized per domain ($z$-scored) and presented as token pairs
(\texttt{param-id}, value); a transformer encoder processes them; two heads
produce the answer (in $\log_{10}$ space for the wide-range beam and RC
domains) and the trajectory (50 points of the dependent variable on a true
grid). The physics-consistency layer then computes how far the predicted
trajectory strays from the governing equation---for the beam, the
moment--curvature relation $EI\,w''=-M(x)$ evaluated on the \emph{true}
grid---and this residual is added to the loss with a warm-up schedule that
lets the data loss settle first.

Two engineering choices proved essential and are, we believe, general
lessons. First, \emph{coordinate columns must not be predicted}: an early
version of the beam model predicted both the position $x$ and the deflection
$w$; because $x$ spans metres while $w$ spans millimetres, the trajectory
loss was dominated by the coordinate, the model fit $x$ and ignored $w$, and
the resulting ``deflection curve'' was physically meaningless (up to
$\sim6\times$ the peak deflection). Restricting the output to the physical
quantity $w(x)$ on the true grid---a coordinate fully determined by the
parameters---fixed the coordinate problem. Second, the
\emph{physics gradient must be scale-free}: raw physics residuals for
fourth-order problems span many orders of magnitude and, unnormalized,
destabilize training (validation error exploded from $7\times$ to
$140\times$ relative error in early experiments). Normalizing the physics
loss by its own detached magnitude each batch---a form of gradient surgery
that makes the physics gradient $\mathcal{O}(1)$ relative to the data
loss---restored stable convergence. Third, \emph{shape and scale must be
decoupled}: beam deflections span five orders of magnitude across the
parameter distribution, so globally normalized trajectory targets were
crushed toward zero for small-deflection problems and the trajectory head
regressed toward an amplitude-weighted average instead of the deformation
shape. We therefore train the trajectory head to predict the
\emph{normalized shape} $\hat w(x)=w(x)/w_{\max}\in[0,1]$, which is the
same universal curve for every simply supported beam (it depends only on
$x/L$), and recover the absolute deflection from the answer head, which
already predicts $w_{\max}$ on a log scale. This decomposition---one head
for shape, one for scale---is the design used throughout.

\subsection{Law-Conditioned Attention: one transformer, six physical laws}\label{sec:lca}

The models above are per-domain: each law gets its own network, and physics
enters only as a loss term. Law-Conditioned Attention (LCA, Figure~\ref{fig:lca};
Methods) makes the \emph{governing equation itself} part of the input and
lets one network serve several laws. Each law is tokenized into a fixed
22-operator vocabulary (\texttt{first\_time\_deriv},
\texttt{second\_space\_deriv}, \texttt{self\_advection},
\texttt{moment\_curvature}, \texttt{energy\_conservation},
\texttt{exponential\_relaxation}, \ldots; Supplementary~Table~S1 lists all
eight), embedded into a law vector, and injected into every attention layer
as a cross-attention stream and into the shared readout as a scale/bias
gate. Input parameters are labeled by \emph{physical quantity} from a shared
13-token vocabulary, so the simply supported beam and the cantilever---which
differ only in support and loading, not in their parameters---present the
\emph{identical} token sequence (\texttt{length}, \texttt{force},
\texttt{modulus}, \texttt{inertia}, \texttt{thickness}). The law signature
is therefore not a task id: it is the only information that distinguishes
two otherwise identical problems, and the model must learn to use it.

\begin{figure}[t]
\centering
\includegraphics[width=0.98\textwidth]{fig/fig12_lca.png}
\caption{\textbf{Law-Conditioned Attention.} The governing equation is
tokenized into a fixed operator vocabulary; the resulting binary signature
is embedded and injected into every transformer layer as a cross-attention
key/value stream, and gates the shared answer/trajectory heads. Input tokens
are physical quantities from a shared vocabulary, so beam and cantilever
present identical inputs and differ only by their law signatures.}
\label{fig:lca}
\end{figure}

We trained one shared transformer (shared body \emph{and} shared heads) on
six laws---beam, cantilever, projectile, pendulum, spring--mass, and RC---at
96 samples per law (576 total, 37\% of the specialist budget), and
compared it against an ablation whose conditioning stream carries no
information (a constant signature; identical architecture, data, and
budget), across six seeds (Figure~\ref{fig:multi}; Table~\ref{tab:lca}).

The cleanest effect of the mechanism is on trajectory fidelity: the
law-conditioned model's held-out curve error (mean absolute deviation from
the true trajectory, normalized by its peak) is $0.093$ versus $0.118$ for
the ablation---a 21\% median reduction, significant by the paired Wilcoxon
test over all 36 (domain, seed) runs ($p=0.0003$, Cliff's
$\delta=+0.33$). Feeding the equation into the transformer makes its
predicted trajectories more physical, even though both conditions receive
the identical per-sample physics loss. This mirrors, at the architectural
level, the residual-level finding of the baselines study: physics
conditioning buys trajectory consistency.

\begin{figure}[t]
\centering
\includegraphics[width=0.98\textwidth]{fig/fig13_multi_law.png}
\caption{\textbf{One transformer, six laws: LCA vs. the constant-signature
ablation} (mean~$\pm$~std over six seeds, held-out answer relative error).
Beam and cantilever present identical parameter tokens; the law signature
is the only signal separating them, and LCA reduces their median error by
19\% and 18\%. On RC, whose normalized trajectory is the same curve for
every time constant (see text), the shared-head generalist fails and the
conditioning stream is significantly harmful---an honest negative result
that localizes \emph{where} law-conditioning pays.}
\label{fig:multi}
\end{figure}

On scalar answers the mechanism's value is concentrated exactly where the
parameters do not identify the law (Table~\ref{tab:lca},
Figure~\ref{fig:multi}). For the two bending laws, which share all five
parameters, LCA reaches 34.0\% and 28.3\% held-out answer error versus
40.4\% and 34.6\% for the ablation---19\% and 18\% median reductions---and
wins on the projectile (5/6 seeds) whose \texttt{speed}/\texttt{angle}
inputs are nearly shared with the pendulum. The pooled answer difference
over all six laws is not significant ($p=0.41$), and one law, RC, is
significantly \emph{worse} with conditioning ($p=0.031$). The RC case is
instructive rather than mysterious: its trajectory is $V(t)=V_0(1-e^{-t/\tau})$,
so after per-channel peak normalization every RC curve is the identical
shape $1-e^{-s}$ for $s\in[0,4]$---the trajectory carries no information
about the answer $\tau=RC$, which must be computed as pure parameter
arithmetic on a log scale that spans six decades. The shared-head
generalist underfits this regime for both conditions, and the conditioning
stream makes it slightly worse. This is an honest negative result: the
mechanism pays where the law disambiguates otherwise-identical inputs, is
neutral where the inputs already identify the law (spring), and can hurt
where the shared head must solve trajectory-independent parameter
arithmetic (RC). We report all six laws rather than dropping the
embarrassing one.

\begin{table}[t]
\centering
\small
\caption{\textbf{One transformer, six laws.} Held-out answer relative error
(median over six seeds); curve error and physics residual are in
\texttt{paper/fig/multi\_law\_data.json}. Beam and cantilever present
identical quantity tokens; the law signature is their only distinguishing
input. Specialists are single-law models at 256 samples (2.7$\times$ more
data per law).}
\label{tab:lca}
\begin{tabular}{@{}lcccc@{}}
\toprule
Domain & LCA & Ablation & Specialist & Ratio vs.\ spec. \\
\midrule
Beam (identical tokens) & 0.340 & 0.404 & 0.103 & 3.3$\times$ \\
Cantilever (identical tokens) & 0.283 & 0.346 & 0.100 & 2.8$\times$ \\
Projectile & 0.110 & 0.112 & 0.092 & 1.2$\times$ \\
Pendulum & 0.021 & 0.023 & 0.008 & 2.6$\times$ \\
Spring & 0.063 & 0.058 & 0.056 & 1.1$\times$ \\
RC (trajectory-degenerate) & 0.611 & 0.509 & 0.360 & 1.7$\times$ \\
\bottomrule
\end{tabular}
\end{table}

Against the per-domain specialists, the generalist is within $1.1$--$2.8\times$
of single-law accuracy on five of six laws at 37\% of the data per law,
and matches the specialist on the two simplest laws (projectile 1.2$\times$,
spring 1.1$\times$). The result we claim is architectural and falsifiable:
one network, six laws, with the equation---not a task id---as the steering
signal, a significant gain in trajectory fidelity, and a precise account
of when the steering matters.

\subsection{Causal steering by the governing equation}\label{sec:steering}

The training contrast (real vs.\ constant signature) shows that the equation
vector is \emph{used}, but not that it \emph{causes} behavior. To test
causality directly we swap the signature at inference: a beam problem---its
tokens and values untouched---is presented to the network with the
cantilever's law signature, and vice versa. Because the two bending laws
share all five parameter quantities and their sampling ranges
(Table~\ref{tab:lca}), the inputs are literally identical; only the injected
equation differs. Any change in the output is therefore caused by the
equation vector alone.

Figure~\ref{fig:steer} shows the result on a representative held-out sample.
Swapping the signature leaves the constant-signature ablation's prediction
exactly untouched (its conditioning stream carries no information), but
displaces the law-conditioned model's prediction substantially, moving it
across the decision boundary between the two bending solutions. Over all
held-out samples ($n{=}72$ per condition, six seeds), the steering
index---positive when the swapped prediction is closer to the target law's
exact solution than to the source law's---is $+0.090$ for LCA versus
$-0.132$ for the ablation (paired Wilcoxon $p<0.0001$), and the
signature-induced disruption (how far the swap moves the prediction away
from the source-law solution) is $0.371$ versus $0.000$ ($p<0.0001$); the
per-seed medians differ with $p=0.031$. The governing equation is thus
causally active at inference: it does not merely label the training
distribution, it steers the network's behavior in the direction of the law
it encodes.

Two qualifications matter. First, steering is \emph{partial}: the swapped
prediction moves toward the target law's solution family but does not land
exactly on it. The model has learned conditional behavior for the pairings
it was trained on, not an explicit equation-to-solution function that
composes at inference. Second, the scalar answer does not re-target under
the swap---the shared head holds its per-domain output scale even when the
signature changes. Both limitations are expected of a mechanism trained
only on matched (tokens, signature) pairs, and they motivate the
adaptation experiment below: the equation cannot re-target behavior
without data, but it should make re-targeting \emph{with} data cheaper,
which is exactly what Section~\ref{sec:fewshot} tests.

\begin{figure}[t]
\centering
\includegraphics[width=0.98\textwidth]{fig/fig14_law_swap.png}
\caption{\textbf{Causal steering by the governing equation.} A held-out beam
problem (identical tokens and values) is presented with the cantilever's law
signature, and a cantilever problem with the beam's, at inference time. The
constant-signature ablation (bottom row) is exactly insensitive---its
prediction does not move. The law-conditioned model (top row) leaves the
source-law solution and moves toward the injected law's exact solution;
pooled steering index $+0.090$ vs.\ $-0.132$ ($p<0.0001$);
signature-induced disruption $0.371$ vs.\ $0.000$ ($p<0.0001$). Steering is
partial: the swapped prediction approaches but does not reach the target
solution.}
\label{fig:steer}
\end{figure}

\subsection{Few-shot adaptation to a new law}\label{sec:fewshot}

The engineer's claim is not only that one network serves the laws it has
seen, but that it generalizes to laws it has not. We held the cantilever
out of pretraining entirely (five laws, 96 samples each), then fine-tuned
the pretrained generalist on 24 cantilever samples---25\% of the
specialist's budget---for 40 epochs, and compared three conditions under
identical data and compute: the LCA generalist (real signature), the
constant-signature ablation, and a from-scratch single-law specialist
trained on the same 24 samples for the same 40 epochs
(Figure~\ref{fig:pvfs}, Table~\ref{tab:fewshot}).

The pretrained generalist adapts to the held-out law with 24 samples far
more accurately than a from-scratch specialist on the same data: median
held-out answer error $0.197$ (LCA) and $0.239$ (ablation) versus $0.580$
for the specialist---$2.9\times$ and $2.4\times$ lower---and curve error
$0.137$ and $0.131$ versus $0.282$ ($2.1\times$ lower). Both generalists
reach the specialist's full-data trajectory fidelity at a quarter of its
data, and the equation signature adds a further 18\% reduction in answer
error ($0.197$ vs.\ $0.239$, two of three seeds; Table~\ref{tab:fewshot}).
The signature's value concentrates precisely where the zero-shot law-swap
could not reach: given a little data, the law vector tells the network
which behavior to route to, so adaptation to a new governing equation costs
a quarter of the specialist's data rather than a full retraining. The
pretrained beam knowledge transfers because beam and cantilever share their
parameter tokens; the equation signature---not the data---tells the adapted
model which deflection family is in force.

\subsection{Physics supervision at the field level}\label{sec:physvdata}

The physics channel acts in two places: the residual loss (all models) and
the law-conditioning input (LCA). The baselines study documented the
residual loss's effect on scalar answers; here we ask whether it improves
held-out fidelity of the full 2D \emph{field}---the regime where a
data-only network must interpolate between sparse training fields. On the
heat plate we trained identical networks with the physics term on
($w_{\mathrm{phys}}=0.05$) and off (data-only), at 25\% and 100\% of the
256-sample budget, three seeds each, 240 epochs
(Figure~\ref{fig:pvfs}, panel a).

On identical budgets the physics term reduces the predicted field's
violation of the governing equation by 4--8$\times$ (mean residual
$0.06$--$0.11$ vs.\ $0.35$--$0.81$) but does not improve---and at this
capacity slightly degrades---held-out fidelity against the exact field
(field max error $0.28$ vs.\ $0.03$ at 256 samples; $0.51$ vs.\ $0.15$ at
64 samples; Figure~\ref{fig:pvfs}a). This is an honest negative that
sharpens the paper's mechanism claim: the \emph{residual-loss} physics
channel buys physical consistency, not field accuracy, when the data
already determine the field; the accuracy benefit of physics must come from
the \emph{input} channel (LCA), which the trajectory-fidelity and
adaptation results above show it does. It also corrects an over-optimistic
earlier claim: the canonical heat-plate case (5.9\% field error) is the
easiest member of the validation distribution, whose mean field error is
$\sim$29\%---a distinction now reported explicitly.

\begin{figure}[t]
\centering
\includegraphics[width=0.98\textwidth]{fig/fig15_physvdata_fewshot.png}
\caption{\textbf{Physics supervision at the field level and few-shot
adaptation.} (a) Held-out 2D field error (max over the plate, mean over
three seeds) on the heat2d domain, identical networks and budgets, with the
physics-consistency loss on ($w_{\mathrm{phys}}=0.05$) or off (data-only),
at 25\% and 100\% of the data. (b) Held-out cantilever answer error after
adapting to the held-out law with 24 samples (25\% of the specialist
budget) for 40 epochs: LCA generalist, constant-signature ablation, and a
from-scratch specialist on the same data.}
\label{fig:pvfs}
\end{figure}

\begin{table}[t]
\centering
\small
\caption{\textbf{Few-shot adaptation to a held-out law.} The cantilever is
absent from pretraining; all conditions see the same 24 samples (25\% of
the specialist budget) for 40 epochs and are evaluated on the same 12
held-out problems. Answer / curve error by seed (median bold).}
\label{tab:fewshot}
\begin{tabular}{@{}lccc@{}}
\toprule
Seed & LCA generalist & Ablation & Specialist \\
\midrule
0 & 0.197 / 0.137 & 0.239 / 0.131 & 0.279 / 0.273 \\
1 & 0.287 / 0.134 & 0.437 / 0.145 & 0.580 / 0.284 \\
2 & 0.182 / 0.172 & 0.153 / 0.129 & 0.691 / 0.282 \\
Median & \textbf{0.197} / 0.137 & 0.239 / 0.131 & 0.580 / 0.282 \\
\bottomrule
\end{tabular}
\end{table}

\subsection{Learning dynamics and generalization}

Figure~\ref{fig:convergence} shows training dynamics for two representative
domains. Validation relative error falls from $54.8\%$ to $4.0\%$
(projectile) and from $13.3\%$ to $2.4\%$ (spring) over 25 epochs, while the
physics residual falls from $5.6$ to $0.05$ and from $361$ to $0.16$,
respectively: the network becomes simultaneously more accurate \emph{and}
more physical, confirming that the data and physics terms cooperate rather
than compete.

\begin{figure}[t]
\centering
\includegraphics[width=0.98\textwidth]{fig/fig3_convergence.png}
\caption{\textbf{Convergence on data and physics together.}
(a) Validation relative error (log scale) for projectile and spring domains.
(b) Physics residual (log scale), i.e.\ deviation from the governing
equation, measured on the validation set; the dashed line marks residual
$0.1$. The physics loss ramps in after epoch 10; both metrics continue to
improve jointly.}
\label{fig:convergence}
\end{figure}

On 250 held-out problems per domain (Figure~\ref{fig:predictions}), the
trained \PhysFormer achieves mean relative errors of $14.8\%$ for simply
supported beams, $11.6\%$ for cantilevers, $9.6\%$ for projectiles, and
$1.9\%$ for the heat plate,
with coefficients of determination $R^2\approx0.99$ (computed in
$\log_{10}$ space for the wide-range bending answers). These are the
\emph{honest} numbers at the 256-sample budget with proper seeding: the
seed-to-seed spread is $\pm0.7$--$2$ percentage points, and the
error concentrates at the small-deflection end of the beam distribution,
where the answer spans five orders of magnitude---precisely the regime
where the physics residual term should help most, and the target of the
trajectory-fidelity work below.

\begin{figure}[t]
\centering
\includegraphics[width=0.98\textwidth]{fig/fig4_predictions.png}
\caption{\textbf{Generalization to held-out problems} (250 per domain,
seeds independent of training). (a) Simply supported beam: predicted vs
exact maximum deflection (log--log; answers span five orders of
magnitude). (b) Cantilever beam: predicted vs exact maximum deflection.
(c) Projectile: predicted vs exact range. (d) Heat plate: predicted vs
exact peak temperature. Dashed lines are identity;
mean relative errors and $R^2$ are given in each panel.}
\label{fig:predictions}
\end{figure}

Figure~\ref{fig:deflection} examines trajectory fidelity for the canonical
beam: the scalar answer is within $8.7\%$ of the closed form and the full
deflection curve within $19\%$ of the peak deflection along the span. The
sigmoid-constrained shape head never produces negative deflections, and
the curve error is a factor of two smaller than the unconstrained head's.
Together with the scalar results of Figure~\ref{fig:predictions}, this
demonstrates \emph{consistent} engineering predictions: the same network
that answers ``what is the maximum deflection?'' also produces the full
deformation shape, and both are verifiable against independent numerics.

\begin{figure}[t]
\centering
\includegraphics[width=0.98\textwidth]{fig/fig5_deflection.png}
\caption{\textbf{Trajectory fidelity for the canonical beam} ($L=4$ m,
$P=3$ kN, $E=2\times10^{11}$ Pa, $I=5\times10^{-6}$ m$^4$).
(a) Deflection curve: exact closed form (solid) vs \PhysFormer (dashed),
with the discrepancy shaded. (b) Normalized error profile along the span.}
\label{fig:deflection}
\end{figure}

\subsection{Ablations: the physics loss and the shape/scale decomposition}

Two design decisions are load-bearing for the physics core, and we ablate
each by retraining the beam model on the same 256 samples for 40 epochs with
a single change (Figure~\ref{fig:ablations}).

\emph{Physics loss on vs off.} Turning off the physics-consistency term
($w_p = 0$) leaves the data loss unchanged. At 40 epochs the scalar answer
accuracy is comparable (validation relative error $10.8\%$ without vs
$13.0\%$ with physics, single seed); at 60 epochs over five seeds the
physics-on model is more accurate (see Table~\ref{tab:baselines}). The
decisive difference is the \emph{trajectory}: the physics-off model's
predicted deflection curve violates the moment--curvature relation by an
order of magnitude (residual $17$ vs $1.8$ at 40 epochs; $\sim$70 vs
$\sim$3 at 60 epochs). The physics term therefore buys what the data loss
cannot provide---consistency with the governing equation---and this
advantage is robust to budget and seed. This matches the cooperation seen
in Figure~\ref{fig:convergence} for the first-order domains, where physics
and data improve together.

\emph{Shape norm vs global norm.} Normalizing every trajectory by its own
peak (shape norm) rather than by dataset-wide statistics is not a cosmetic
choice: with global normalization the trajectory targets are crushed toward
zero for the small-deflection problems that dominate the distribution, the
trajectory head regresses toward an amplitude-weighted average, and---because
the encoder is shared---the answer head degrades too (validation relative
error $59\%$ vs $13\%$, and the predicted curve violates the
moment--curvature relation by four orders of magnitude). With shape norm,
the head learns the universal deformation shape and scales it by the answer
head, and the trajectory is physical. The global-norm curve is not an
engineering prediction at all.

\emph{Negative results on trajectory capacity.} Because shape fidelity was
the remaining gap, we also tried to strengthen the trajectory head and its
supervision: a three-layer, wider trajectory MLP, and a hard-example physics
weighting that emphasizes each batch's worst physical violations. Both
\emph{hurt} (validation error $\sim$13--16\% vs $\sim$6--14\% for the
released configuration on the same budget), because the larger head's
trajectory gradients perturb the shared encoder and the residual weighting
amplifies outlier gradients. These negative results are reproducible from
the committed training logs and code; they bound the sensible capacity at
this data scale.

\begin{figure}[t]
\centering
\includegraphics[width=0.98\textwidth]{fig/fig6_ablations.png}
\caption{\textbf{Ablations (beam, 40 epochs, 256 samples, same architecture).}
(a) Validation relative error with and without the physics loss: comparable
scalar accuracy, while the physics-on trajectory residual is an order of
magnitude lower (1.8 vs 17). (b) Physics residual of the predicted
trajectory under per-problem shape normalization vs dataset-wide (global)
normalization: the global-norm trajectory is unphysical by four orders of
magnitude, and its answer error is $\sim$4.5$\times$ higher (59\% vs 13\%).}
\label{fig:ablations}
\end{figure}

\subsection{The cantilever domain: does the design transfer?}

The simply supported beam is a special case: its normalized deflection
shape depends only on $x/L$, which makes the shape/scale decomposition
almost too easy. The cantilever beam under a tip load is the natural
stress test: a different boundary condition (clamped at one end, free at
the other) and a threefold larger deflection for the same load
($w_{\max}=PL^3/3EI$ vs $PL^3/48EI$), but the same moment--curvature
physics. We added it as a sixth domain with no new machinery: the
closed-form trajectory, an independent finite-difference verifier that
integrates the static moment distribution with a ghost point at the clamped
end, and a training residual identical in form to the beam's. The trained
\PhysFormer (same architecture and budget as the beam) reaches held-out
accuracy of $10.4\%$ (Figure~\ref{fig:cantilever}), and its deflection
curve for the canonical cantilever is within $3\%$ on the scalar and $16\%$
of the peak along the span---the shape/scale decomposition transfers
without modification. This is
evidence that the physics layer is a \emph{mechanism}---a constraint that
adapts to any domain carrying a governing equation---rather than a
hand-tuned decoder for one problem.

\begin{figure}[t]
\centering
\includegraphics[width=0.98\textwidth]{fig/fig7_cantilever.png}
\caption{\textbf{Cantilever beam under a tip load} ($L=4$ m, $P=3$ kN,
$E=2\times10^{11}$ Pa, $I=5\times10^{-6}$ m$^4$). (a) Deflection curve:
exact closed form vs \PhysFormer, discrepancy shaded. (b) Held-out
predictions (250 problems): predicted vs exact maximum deflection.}
\label{fig:cantilever}
\end{figure}

\subsection{Baselines: attention and physics, isolated and measured}

All of the above used a single configuration. To place it against
alternatives, we trained three model kinds on identical budgets (256 samples,
60 epochs, the same answer/trajectory heads, five seeds each, seeds
controlling data, shuffling, and weight initialization): an MLP encoder with
the same capacity (no attention, no physics), the transformer without the
physics term, and \PhysFormer with the physics term
(Table~\ref{tab:baselines}, Figure~\ref{fig:baselines}).

Two findings stand out. First, the physics term's benefit is most robust at
the \emph{trajectory} level: the predicted deflection curves of the
physics-on models violate the governing equations by one order of magnitude
less than the no-physics transformer on both bending domains (beam $3.4$
vs $48$; cantilever $39$ vs $556$), by an order of magnitude less on
Burgers ($0.01$ vs $0.09$), and by more than an order of magnitude less on
the heat plate ($0.45$ vs $10.4$). The physics loss is the only mechanism
that penalizes consistency with the equations directly, and it does so
across all domains. Second, on scalar accuracy the picture is domain-dependent:
\PhysFormer matches or beats the MLP on the bending and conservation-law
domains (beam $14.7\%$ vs $15.6\%$; cantilever $14.8\%$ vs $15.2\%$;
Burgers $0.9\%$ vs $1.8\%$) and the no-physics transformer everywhere
(which also shows the noisest seed-to-seed spread, e.g.\ beam $19.3\%\pm6.3$
vs $14.7\%\pm0.7$), but on the simplest two-parameter domain (projectile)
the MLP wins ($6.5\%$ vs $10.8\%$): with only two parameters and a nearly
linear map, attention and physics are not needed. This is an honest
negative---the extra machinery pays for itself only when the
parameter-to-answer map and its trajectory are rich enough.

\emph{Are these differences significant?} With five seeds per kind, the
per-domain Wilcoxon signed-rank test is limited to $p\geq0.0625$ (the
finest resolution of a five-pair sign permutation), so we pool all five
domains into 25 paired runs for each comparison (the script
\texttt{physx/significance.py} reproduces every number from the training
logs). On the trajectory residual, pooling gives $p<0.001$ for
physics-on vs no-physics with Cliff's $\delta=-0.92$ (a \emph{large}
effect, in the direction of physics-on) and a $\sim$19-fold median
reduction; physics-on vs MLP gives $p<0.001$, $\delta=-0.84$ (3-fold
median reduction); and the no-physics transformer is significantly
\emph{worse} than the MLP ($p<0.001$, $\delta=+0.84$), consistent with the
unconstrained attention head's wider variance. On the scalar answer,
pooling shows \emph{no} significant difference between any pair of kinds
($p=0.65$, $0.79$, $0.86$ respectively), with effect sizes
$|\delta|\leq0.12$. The physics term therefore does not move the scalar
answer---it buys physical consistency of the full trajectory, and it does
so significantly and with large effect. (Per-domain $p$-values and
complete data are in \texttt{paper/fig/significance.json}.)

\begin{table}[t]
\centering
\small
\begin{tabular}{lcccc}
\toprule
 & \multicolumn{2}{c}{Validation rel.\ error (\%)} & \multicolumn{2}{c}{Trajectory residual} \\
\cmidrule(lr){2-3}\cmidrule(lr){4-5}
Domain & MLP & Transformer & $-$physics & $+$physics \\
\midrule
Beam & $15.6\pm2.3$ & $19.3\pm7.0$ & $48\pm14$ & $3.4\pm0.5$ \\
Cantilever & $15.2\pm1.5$ & $20.4\pm7.5$ & $556\pm214$ & $39\pm12$ \\
Projectile & $6.5\pm1.9$ & $7.6\pm5.6$ & $0.12$ & $0.10$ \\
Burgers & $1.8\pm0.7$ & $1.0\pm0.3$ & $0.09$ & $0.01$ \\
Heat2D & $1.6\pm0.4$ & $1.4\pm0.5$ & $10.4\pm5.9$ & $0.45\pm0.02$ \\
\bottomrule
\end{tabular}
\caption{\textbf{Equal-budget baselines} (256 samples, 60 epochs, five seeds,
mean~$\pm$~std). Left: validation relative error of the scalar answer for
the MLP and no-physics transformer; the right-hand columns report the
physics-on model's trajectory residual, the mean governing-equation
violation of the predicted trajectory on the validation set (the MLP and
no-physics values are omitted for space; the MLP residuals are $4.6$,
$62.6$, $0.12$, $0.06$, $1.33$ for the five rows---always higher than
physics-on except for the trivial projectile where all three are small).
The physics term consistently reduces the trajectory's violation of the
governing equation by about an order of magnitude; pooling all 25 paired
runs, the effect is significant (Wilcoxon $p<0.001$, Cliff's
$\delta=-0.92$; see text).}
\label{tab:baselines}
\end{table}

\begin{figure}[t]
\centering
\includegraphics[width=0.98\textwidth]{fig/fig9_baselines.png}
\caption{\textbf{Multi-seed baselines} (mean~$\pm$~std over five seeds,
identical budgets). The physics term's most robust benefit is trajectory
consistency (one order of magnitude lower governing-equation residual on
beam, cantilever, and Burgers); on scalar accuracy it matches or beats the
MLP where the parameter space is rich, and the no-physics transformer
shows the widest seed-to-seed spread.}
\label{fig:baselines}
\end{figure}

\subsection{The Burgers equation: a nonlinear conservation law}

The seven preceding domains are all linear or weakly nonlinear. The viscous
Burgers equation $u_t + u u_x = \nu u_{xx}$ is a different regime: a
nonlinear conservation law whose solutions steepen into shocks as $\nu$
decreases (Figure~\ref{fig:burgers}a). It stresses every design choice: the
predicted object is a full spatiotemporal field $u(x,t)$ on a $25\times8$
grid (200 trajectory points, four times the bending domains), the exact
reference comes from the Cole--Hopf transform (evaluated by direct heat-
kernel quadrature, since an FFT derivative rings on the non-periodic,
many-orders-of-magnitude $\phi_0$), and the independent verifier is a
finite-volume upwind solver written from the conservation form
(Table~\ref{tab:verify}).

The physics-consistency layer evaluates the Burgers residual pointwise on
the true grid. During training it pulls the predicted field toward
satisfying the equation: the physics-on models reach a trajectory residual
of $0.01$, an order of magnitude below the no-physics transformer
($0.09$) and below the MLP ($0.06$) (Figure~\ref{fig:burgers}c). The scalar
answer---the peak velocity of the final field---is learned almost exactly
(held-out relative error $\sim$1\%; canonical case 0.5\%, shock-forming
case 5\%), and the predicted final-time field resolves the developing
shock profile to $\sim$11\% of the peak for the canonical case
(Figure~\ref{fig:burgers}a); the sharp-gradient shock case remains coarse
on the $25\times8$ grid ($\sim$52\%, see Limitations). The shape/scale
decomposition carries over: the field peaks at the initial amplitude $A$,
so the shape head predicts $u(x,t)/A\in[0,1]$ (constrained by a sigmoid)
and the absolute field is recovered by scaling with the input parameter
$A$.

\begin{figure}[t]
\centering
\includegraphics[width=0.98\textwidth]{fig/fig8_burgers.png}
\caption{\textbf{Viscous Burgers equation} (nonlinear conservation law,
shock formation). (a) Final-time field $u(x,t{=}0.4)$: exact (Cole--Hopf)
vs \PhysFormer for a canonical and a shock-forming case. (b) Held-out
predictions of the peak velocity (250 problems). (c) Governing-equation
residual of the predicted field during training: the physics term drives
it an order of magnitude lower than the no-physics transformer.}
\label{fig:burgers}
\end{figure}

\subsection{The heat plate: a two-dimensional field domain}

Every domain above lives on a one-dimensional trajectory. The eighth
domain, steady-state heat conduction on a square plate, is the first
\emph{two-dimensional} field: the predicted object is $T(x,y)$ on a
$20\times20$ grid (400 points), governed by Poisson's equation
$\nabla^2 T = f$ with the manufactured source
$f=-(A\pi^2(k^2+l^2)/2)\sin(k\pi x)\sin(l\pi y)$ and Dirichlet edges
$T=A/2$; the exact solution is $T=A(1+\sin(k\pi x)\sin(l\pi y))/2$, so
mode numbers $k,l$ and amplitude $A$ parameterize the problem family. The
shape/scale decomposition transfers with only one change: because the peak
$A$ is attained where the sinusoids coincide, which may fall between
sampling nodes, the field is normalized by the \emph{parameter} $A$ rather
than its sampled maximum (the same choice Burgers made). The independent
verifier is a sparse finite-difference Poisson solve written from the PDE
and the source term alone (agreement $<2\times10^{-3}$, Table~\ref{tab:verify});
the physics-consistency layer penalizes the fourth-order-accurate
Laplacian residual $\nabla^2 T - f$ on the interior of the true grid.

Figure~\ref{fig:heat2d} shows the result. Held-out peak-temperature error
is $1.9\%$ ($R^2=0.998$), and the canonical plate's field is reproduced to
$\sim$6\% of the peak (mean error $\sim$1.3\%). The 400-point oscillatory
field needs a longer trajectory budget than the one-dimensional domains:
the 60-epoch model underfits it badly (max error $\sim$50\%---the head
regresses toward the field's mean), and doubling the budget to 240 epochs
with the heavier trajectory weight captures the mode structure (see
Limitations). The trajectory residual
drops from $10.4$ (no-physics) to $0.45$ (physics-on)---the largest
relative improvement of any domain, as the Laplacian constraint couples
every interior point to its neighbors and cannot be satisfied by chance on
a 400-point field. This is the regime where the physics channel pays most
---for \emph{physical consistency}. A controlled physics-vs-data-only
sweep (Section~\ref{sec:physvdata}) shows the residual term cuts the
field's governing-equation violation 4--8$\times$ at equal data without
improving held-out fidelity against the exact field: the canonical
$\sim$6\% figure is the easiest member of the validation distribution,
whose mean field error is $\sim$29\%. Field accuracy is carried by the
data and, at the architectural level, by the LCA input channel, not by the
residual loss.

\begin{figure}[t]
\centering
\includegraphics[width=0.98\textwidth]{fig/fig10_heat2d.png}
\caption{\textbf{Steady-state heat conduction on a square plate} (Poisson
$\nabla^2 T=f$, modes $k{=}2$, $l{=}3$). (a) Exact temperature field
$T(x,y)$; (b) \PhysFormer prediction; (c) pointwise error. Bottom:
held-out peak temperatures (250 problems, rel.\ error $1.9\%$).}
\label{fig:heat2d}
\end{figure}

\subsection{Comparison against an external PINN implementation}

The baselines above are all internal to our codebase. To check against the
standard tooling of the physics-informed learning literature, we trained
\textsc{DeepXDE}\pn{\citep{lu2021}}---the most widely used PINN
library---on the same viscous Burgers problem, giving it the setup most
favourable to a PINN: \emph{one fresh network per problem instance}
(a 2--50$\times$4--1 tanh FNN, 2{,}000 collocation points, Adam + L-BFGS,
12k + 2k iterations, identical for every instance, exact Cole--Hopf
boundary data at $x=\pm1$).
\PhysFormer, by contrast, is a single network trained once on 256 problem
instances that generalizes to unseen $(\nu,A,\sigma)$; the comparison is
therefore conservative against \PhysFormer. Both are evaluated with
identical metrics on the model grid: scalar (final-time peak) error, full
field error, and the same finite-difference governing-equation residual
(Figure~\ref{fig:deepxde}).

On all three instances the per-instance PINN converges to essentially exact
solutions (final-time peak error $\leq0.01\%$; full-field error
$0.03$--$0.06\%$; $\sim$19--21 minutes of training per instance), which
validates both the problem setup and the Cole--Hopf ground truth: a
standard PINN implementation reproduces the reference solutions to near
machine precision given per-instance training and exact boundary data. The
generalist \PhysFormer (one network, one $\sim$minute training, then a
millisecond forward pass for \emph{any} instance) achieves final-time peak
errors of $0.02$--$3.7\%$ across the same instances and full-field errors
of $33$--$51\%$---worse on any given instance, as expected, but with a
qualitatively different cost profile: no per-instance optimization at all.
Notably, \PhysFormer's governing-equation residual on the predicted field
is comparable to or lower than the converged PINN's ($2\times10^{-4}$ vs
$1$--$5\times10^{-3}$), i.e.\ the generalist's predictions are as
physical as the specialist's, just less accurate. The comparison thus
quantifies the generalist--specialist tradeoff directly: the physics-
consistency layer makes a single network physical everywhere, while a
per-instance PINN remains the reference for accuracy on a single problem.

\begin{figure}[t]
\centering
\includegraphics[width=0.98\textwidth]{fig/fig11_deepxde.png}
\caption{\textbf{External PINN baseline (\textsc{DeepXDE}) on viscous
Burgers.} Three representative instances (canonical $\nu{=}0.05$, mild
$\nu{=}0.10$, shock-forming $\nu{=}0.02$); per-instance PINN (one fresh
FNN per problem, $\sim$20 min training each) vs the generalist
\PhysFormer (one network, one forward pass). Left: final-time peak error.
Centre: full $(x,t)$-field error. Right: governing-equation residual of the
predicted field on the model grid (log scale).}
\label{fig:deepxde}
\end{figure}

\subsection{The verification gate: a controlled benchmark}

To measure the gate rather than merely assert it, we ran seven deterministic
missions---benign software tasks, two beam/cantilever/Burgers design tasks,
an injected-fault task, a skill-level guard, and a TODO scan---twice each:
with the gate enforced and with it disabled (\texttt{enforceGate: false},
which makes gated failures indistinguishable from successes, the
verifier-free baseline). Ground truth was computed independently of the
agent by running the test suites ourselves and by the solver's own numeric
verification (Table~\ref{tab:gate}).

On benign missions both arms behave identically. The discriminator is the
fault task: with the gate, the agent reports failure and journals a lesson;
without it, the agent reports success on code whose tests actually fail.
Across all seven missions the gate produces \emph{zero} false successes
while the verifier-free agent claims success on two of seven (the injected
fault and the clobbered-directory guard, whose protection is skill-level
and therefore gate-independent). The verified-success rate is unchanged
(5/7), but the false-success rate drops from 29\% to 0\%. The gate does not
make the agent smarter; it makes its claims \emph{honest}, which is the
property that lets trial-and-error learning be trusted at all.

\begin{table}[t]
\centering
\small
\begin{tabular}{lccc}
\toprule
 & Reported success & Ground truth OK & False success \\
\midrule
With gate & 5/7 & 5/7 & 0/7 (0\%) \\
Without gate & 7/7 & 5/7 & 2/7 (29\%) \\
\bottomrule
\end{tabular}
\caption{\textbf{Verification-gate benchmark} (7 missions, mechanical brain).
Ground truth is computed independently of the agent. The gate turns claimed
success into verified success; the verifier-free baseline claims success on
the two missions whose ground truth is broken (an injected test-suite fault
and a skill-level clobber guard).}
\label{tab:gate}
\end{table}

\subsection{End-to-end engineering mission}

Finally, we exercised the complete agent---not the components---on a physics
design task expressed in natural language:

\begin{quote}
\emph{``design a simply supported beam with span 4 m, point load 3000 N,
modulus 2e11 Pa, inertia 5e-6 m\textsuperscript{4}, height 0.2 m''}
\end{quote}

The mechanical brain recognized the beam domain, parsed all five parameters
from natural language, and emitted a single-step plan invoking the physics
skill. The skill computed the closed-form answer ($4.0$ mm), verified it
against the independent finite-difference solver (relative residual
$2.6\times10^{-9}$), reported the associated maximum stress
($120$ MPa at the given section), and produced a \PhysFormer prediction of
$3.65$ mm---within $8.7\%$ of the exact answer, together with the full
deflection curve (maximum error $19\%$ of the peak deflection; see
Figure~\ref{fig:deflection}). An identical mission for a cantilever
(``design a cantilever beam with length 4 m, tip load 3000 N, modulus
$2\times10^{11}$ Pa, inertia $5\times10^{-6}$ m$^4$, height 0.2 m'') is
recognized, parsed, solved ($64$ mm), and verified ($3.1\times10^{-6}$)
through the same single-step plan, with a \PhysFormer prediction within
$3\%$. A third mission---``design a viscous burgers flow with viscosity
0.05, amplitude 1.5, width 0.3''---is recognized, parsed, solved, verified
against the finite-volume solver, and predicted to within $0.5\%$ of the
exact peak velocity, again through the same single-step plan: adding each
new domain required no changes to the agent loop. Each episode, its
verification outcome, and the model's prediction were journaled. The whole
mission completed in one iteration in under two seconds of solver time.
The same loop, with the LLM brain substituted for the mechanical one,
accepts arbitrary software-engineering goals (``add a feature and make the
tests pass'').

\section{Discussion}

\Age demonstrates that the architecture of an ``artificial general
engineer''---the concept Prometheus is commercializing at a \$41B
valuation\pn{\citep{techcrunch2026}}---can be built today in a few thousand
lines of dependency-light code, with no proprietary infrastructure, and with
every claim checked by execution rather than assertion. Three design
principles emerge that we expect to transfer to larger systems.

First, \emph{verification is the load-bearing principle}. The agent's
authority is not the model but the environment: tests for software,
numerically independent simulators for physics. This mirrors the shift to
process-supervised verifiers in mathematical reasoning\pn{\citep{lightman2024}}
and test-gated code generation\pn{\citep{li2022,jimenez2024}}, and it is the
property that makes trial-and-error learning safe enough to leave the
sandbox. Our injection experiment---silently breaking a test suite and
observing the gate fail the mission and journal a lesson---is a direct
demonstration that the loop's verdicts track ground truth, not the model's
confidence.

Second, \emph{physics should enter the transformer not only as a loss but as
input}. The joint training of reasoning layers and physics-consistency
layers (a PINN-style residual\pn{\citep{raissi2019}}) produced models that
improve on both accuracy and physical consistency simultaneously
(Figure~\ref{fig:convergence}), and the multi-seed baseline study isolates
that contribution: the physics term consistently reduces the predicted
trajectory's violation of the governing equation by about an order of
magnitude, a benefit that the data loss alone cannot provide
(Table~\ref{tab:baselines}). But a loss term is a weak channel: it says
nothing about \emph{which} law governs a problem until after the output
exists. Law-Conditioned Attention (Section~\ref{sec:lca}) closes that gap
by making the governing equation itself part of the input---tokenized into
a fixed operator vocabulary and injected into every attention layer and
into the readout---so a single shared transformer can serve multiple laws
and the model is told, up front, what physics it is being asked to reason
about. The causal test of Section~\ref{sec:steering} confirms the channel
is not decorative: on inputs that are literally identical across two laws,
swapping the injected signature measurably steers the prediction toward
the target law's solution, and a constant-signature control is exactly
insensitive---the equation is doing causal work, not labeling data. And
Section~\ref{sec:fewshot} shows the same channel makes adaptation to a law
held out of pretraining data-efficient, the property a ``general
engineer'' actually needs. Three failures we diagnosed and fixed are, we
suspect, common causes of unstable physics-informed training in the wider
literature\pn{\citep{wang2021}}: predicting redundant coordinate columns;
letting raw residual magnitudes dominate the gradient; and letting the
unconstrained trajectory head emit unphysical values (negative
deflections/velocities) that the data loss is too weak to remove. Our fixes
(normalized true-grid targets; scale-free residual gradients; a sigmoid
constraint on shape-normalized outputs) are cheap and general.

Third, \emph{memory should be episodic and replayable}. The journal makes
every failed verification a lesson available to future plans, converting
individual missions into cumulative experience without weight updates. This
is a deliberately simple answer to the consolidation problem that continual
learning research addresses with more elaborate
mechanisms\pn{\citep{kirkpatrick2017}}.

\subsection{Limitations}

We are explicit about what \Age does \emph{not} yet do. (1) \emph{Scale}: the
physics models are trained on CPU in minutes with 3--7 parameters per domain
and a 48-dimensional transformer; generalization errors of $\sim$10--15\%
at the 256-sample budget are an order of magnitude above what larger budgets
achieve (the 1,500-sample projectile run reached $\sim$4\%), and the beam
error concentrates at the small-deflection tail, where the answer spans
five orders of magnitude. (2) \emph{Trajectory fidelity}: the shape head
resolves the smooth deflection and velocity profiles well (curve errors
$\sim$11--19\% of peak on the showcase cases) but remains coarse where the
field is genuinely hard---the Burgers shock case, whose sharp gradient is
under-resolved on the $25\times8$ grid, shows a $\sim$52\% curve error even
though the scalar is within 5\%. Finer grids, more trajectory capacity, and
residual-weighted sampling are the natural next steps. (3) \emph{Repair}:
the mechanical
brain inspects and solves but does not \emph{author fixes}; closing the loop
on red tests requires the LLM brain and a capable coder model---the loop
supports it, but we have not tuned it end-to-end. (4) \emph{Simulated
physics}: verification is against numeric integration of the same governing
equations, not against sensors, prototypes, or manufactured parts; the
``real-world trial and error'' that is Prometheus's stated thesis is the
next step, and the architecture ports unchanged (swap the verify skill for a
measurement harness). (5) \emph{Physics task classes}: quasi-static beam
bending and scalar conservation laws, not general continuum mechanics or
multi-physics coupling. (6)
\emph{Safety}: confinement to a working directory and a command denylist are
necessary but not sufficient for autonomous operation; we recommend
containerization and human review before any deployment with write access.

\subsection{Implications}

The significance of this work is not the individual components---agents,
PINNs, and journals each have rich literatures---but their integration into
a minimal, testable, honest whole whose failure modes are visible. As
foundation models continue to improve\pn{\citep{brown2020,bubeck2023,bommasani2021}},
the planner will become stronger; the loop, the gate, and the physics core
will remain. The specific next milestones are: (i) verifier-driven repair, in
which the failing diff is fed back to the planner until the gate goes green;
(ii) scaling Law-Conditioned Attention to harder regimes---the causal
steering and few-shot results show the equation channel both steers
behavior and accelerates adaptation, so the natural extension is more laws
in one network, a shared trajectory head over mixed geometries, and larger
operator vocabularies that can express coupled multi-physics problems the
model has never seen; (iii) trajectory-level physics consistency for
fourth-order problems, likely via residual-weighted sampling and
higher-capacity trajectory heads; (iv) coupling to physical simulators
(finite-element, multibody) so verification approaches manufactured
reality; and (v) journal consolidation, so lessons generalize across
domains rather than accumulating. Each milestone is
directly testable in the current framework, and each moves \Age---and the
``general engineer'' concept it instantiates---closer to the physical world
its commercial namesake targets.

\section*{Methods}

\subsection*{Agent loop}

\textbf{Algorithm 1:} \Age trial-and-error loop.
\begin{enumerate}[label=\arabic*., leftmargin=1.6em]
\label{alg:loop}
\item \textbf{Plan.} Brain maps goal $\to$ ordered steps
      $\{(skill, args, expect)\}$; ``expect: ok'' marks gates.
\item \textbf{Act.} For each step, invoke the skill in the sandboxed working
      directory; record $\{ok, status, output\}$.
\item \textbf{Gate.} If a gated step fails, mark mission failed, derive a
      lesson from the output tail, journal, and stop.
\item \textbf{Reflect.} Brain summarizes the iteration; if not done, emits
      the next plan (LLM brain re-plans from all observations).
\item \textbf{Remember.} Append the episode $\{goal, results, summary,
      status, lessons\}$ to the JSONL journal; future plans receive recent
      lessons as context.
\end{enumerate}

The loop runs at most $I$ iterations (default 6); reaching $I$ without a
terminal reflection is itself a failure. The mechanical brain is a
deterministic recognizer of mission classes (scaffold/explain/scan/test/
physics) with hand-written natural-language parameter extraction (e.g.\
\texttt{span 4 m} $\to$ $L=4$; \texttt{modulus 2e11 Pa} $\to$ $E=2\times10^{11}$);
the LLM brain emits JSON plans through any OpenAI-compatible chat endpoint
and degrades gracefully to inspect-only on error.

\subsection*{Embodiment and sandboxing}

All skills operate relative to a root directory. Path operations are guarded
by a containment check (resolved target must stay within the root).
Commands are executed with a 120\,s timeout and a denylist that blocks
destructive patterns (\texttt{rm -rf /}, \texttt{git push},
\texttt{git reset --hard}, privilege escalation, disk formatting, and
shell-pipe-to-\texttt{sh}) unless explicitly overridden by an environment
flag. A leading \texttt{python} token is resolved to an interpreter present
on the host. Python bytecode caches (\texttt{\_\_pycache\_\_}) under the
working directory are purged before every python command: a source edit that
is same-length and lands within the same timestamp second as a cached
compile is invisible to Python's pyc invalidation check, and an uncached
verify step would run stale code and pass broken sources (we observed
exactly this failure mode and made the purge part of the gate). The physics skill spawns the solver without a shell, passing
JSON via argv to avoid quoting hazards across platforms.

\subsection*{Physics domains and independent verification}

For each domain the simulator module provides \texttt{closed} (exact
solution), \texttt{trajectory} (reference curve), and \texttt{verify}
(independent numerics). Governing equations and verifiers:

\begin{itemize}[leftmargin=1.6em,itemsep=1pt]
\item \textbf{Projectile:} $y = x\tan\alpha - gx^2/(2v_0^2\cos^2\alpha)$;
      verifier is Euler integration at $\Delta t = T/2000$; tolerance $2\times10^{-3}$.
\item \textbf{Pendulum:} $\ddot\theta = -(g/L)\sin\theta$; verifier is RK4 at
      $\Delta t=10^{-4}$ (period convergence under step halving); tolerance $10^{-3}$.
\item \textbf{Spring--mass:} $m\ddot x + kx = 0$; verifier is the ODE
      residual by central differences; tolerance $10^{-3}$.
\item \textbf{Beam:} $EI w'''' = q$ with simply supported boundary
      conditions; closed form $w_{\max} = PL^3/(48EI)$ and deflection
      $w(x)$; verifier is a 400-node finite-difference solver with ghost
      points for the moment-free ends; tolerance $2\times10^{-2}$. The
      training-time physics residual uses the moment--curvature relation
      $EI w'' = -M(x)$, which holds pointwise for a central point load
      (unlike the fourth-order ODE, which carries a Dirac source).
\item \textbf{Cantilever:} $EI w'' = -M(x)$, $M(x) = -P(L-x)$ for a tip
      load; closed form $w_{\max} = PL^3/(3EI)$ and deflection
      $w(x) = P x^2(3L-x)/(6EI)$; verifier is a 400-node finite-difference
      double integration of the static moment distribution with a ghost
      point at the clamped end ($w_{-1} = w_1$ for the zero slope);
      tolerance $2\times10^{-2}$.
\item \textbf{Burgers:} $u_t + u u_x = \nu u_{xx}$ on $x\in[-1,1]$
      with $u(x,0) = A\exp(-x^2/2\sigma^2)$, $t\in[0,0.4]$; the exact
      reference comes from the Cole--Hopf transform
      ($u = -2\nu\,\phi_x/\phi$ with $\phi$ the heat evolution of
      $\phi_0(x) = \exp(-(1/2\nu)\int_{-\infty}^{x}u_0)$), evaluated by
      direct kernel quadrature on a fine grid---an FFT derivative rings on
      the non-periodic, many-orders-of-magnitude $\phi_0$; the predicted
      trajectory is the flattened $25\times8$ (x, t) field on the true
      grid; the verifier is a 2{,}000-cell finite-volume upwind solver for
      the conservation form $u_t + (u^2/2)_x = \nu u_{xx}$ (no reference
      to the closed form); tolerance $5\times10^{-2}$. The scalar answer is
      the peak of the final field, and the field is reconstructed as
      shape $\times$ $A$ at inference (the field peaks at the initial
      amplitude $A$).
\item \textbf{RC circuit:} $\dot V + V/\tau = V_0/\tau$, $\tau=RC$; verifier
      is Euler integration at $\Delta t = 4\tau/4000$; tolerance $2\times10^{-3}$.
\item \textbf{Heat2D:} Poisson $\nabla^2 T = f$ on the unit square with
      $f = -(A\pi^2(k^2+l^2)/2)\sin(k\pi x)\sin(l\pi y)$ and Dirichlet
      edges $T=A/2$; closed form $T = A(1+\sin(k\pi x)\sin(l\pi y))/2$
      (peak $A$); verifier is a sparse finite-difference Poisson solve on a
      $48\times48$ grid written from the PDE and source term alone
      (tolerance $2\times10^{-3}$); the predicted trajectory is the
      flattened $20\times20$ field, and the field is reconstructed as
      shape $\times$ $A$ at inference (the peak $A$ need not fall on a
      sampling node).
\end{itemize}

\subsection*{Dataset and normalization}

Problems sample uniform parameter ranges per domain (e.g.\ beam:
$L\in[1,6]$ m, $P\in[10^3,5\times10^4]$ N, $E\in[69,210]$ GPa,
$I\in[10^{-7},2\times10^{-5}]$ m$^4$, section height $h\in[0.05,0.4]$ m;
Burgers: $\nu\in[0.02,0.1]$, $A\in[0.6,2]$, $\sigma\in[0.2,0.5]$).
Answers spanning orders of magnitude (beam and cantilever deflection, RC
time constant) are learned in $\log_{10}$ space. Parameters are $z$-scored
per domain. Trajectories for the first-order domains are standardized by
global per-channel mean/std, where the wide dynamic range is contained; for
the bending domains, Burgers, and RC, whose dependent variables span
several orders of magnitude across the parameter distribution, each
trajectory is normalized by its own peak (\texttt{shape-norm}), so the
head learns the normalized shape $\hat w(x)=w(x)/w_{\max}\in[0,1]$
(universal for the simply supported and cantilever beams, parameter-
dependent for Burgers and heat2d) and the absolute curve is recovered by
scaling with the appropriate scale (the answer head's prediction for the
bending domains, $V_0$ for RC, and the amplitude parameter $A$ for Burgers
and heat2d). For
shape-norm domains the trajectory head output is passed through a sigmoid,
which constrains the predicted shape to $(0,1)$---deflections, voltages,
velocities, and temperatures in these domains are non-negative, and the
constraint removes unphysical negative predictions. Each problem stores parameters,
answer, unit, and the reference trajectory.

\subsection*{\PhysFormer architecture and training}

Input tokens are (param-id embedding $+$ normalized-value projection $+$
positional encoding). Reasoning layers are a standard transformer encoder
($d_{\text{model}}=48$, $4$ heads, 3 layers, GELU, dropout $0.05$); the
MLP baseline replaces the encoder with an equal-capacity MLP over the
flattened tokens and is trained data-only. The answer head is two linear
layers; the trajectory head is an MLP with a single hidden layer of width
64 (recorded in the checkpoint) projecting the final hidden state to
$T \times C$ points where $T=50$ for the one-dimensional domains,
$T=200$ for the Burgers field, and $T=400$ for the heat plate, and $C=2$
for projectile/pendulum/spring and
$C=1$ otherwise (single dependent column on a true grid). The trajectory
capacity is deliberate: reproducing a full deformation curve is a harder
task than regressing one scalar, and the ablations hold it fixed when
isolating the physics loss.
Training minimizes

\[
\mathcal{L} = \underbrace{\|a - a^*\|^2}_{\text{answer}} \;+\;
\underbrace{w_t\,\|T - T^*\|^2}_{\text{trajectory data}} \;+\;
w_p(t)\,\underbrace{\frac{\rho}{\rho_{\text{det}} + \epsilon}}_{\text{scale-free physics}},
\]

with $w_t=1$ for all domains (the Burgers headline model uses $w_t=3$ and
120 epochs because its field is four times larger).

where $\rho$ is the mean governing-equation residual of the predicted
(de-scaled) trajectory, $\rho_{\text{det}}$ is its detached magnitude
(rendering the physics gradient scale-free), and the weight
$w_p(t) = 0.05\min(1, t/10)$ ramps the physics term in after the data loss
settles. We also explored \emph{hard-example weighting} of the physics
term (each sample's residual weighted by its own detached magnitude before
averaging, a residual-adaptive refinement in the spirit of PINN hard
sampling) and found that at this data scale it amplifies outlier gradients
and degrades both accuracy and consistency (validation error $\sim$14\%
vs 6--13\% without it), so the released models use the plain scale-free
form; the option remains available in the training code. Optimization is
AdamW ($\mathrm{lr}=10^{-3}$, weight decay $10^{-4}$) with global gradient
norm clipping at $1.0$. Bending/RC
answers are inverse-transformed from $\log_{10}$ for evaluation. The seed
controls data generation, batch shuffling, \emph{and} weight initialization
($\texttt{torch.manual\_seed}$), so multi-seed runs are proper statistical
replicates; the baselines table reports mean~$\pm$~std over five seeds.
All headline models are trained on 256 samples per domain at 60 epochs on a
CPU (seconds to minutes per run; the earlier 1,500-sample projectile run
reached $\sim$4\% error, showing that larger budgets monotonically
improve accuracy).

\subsection*{Law-Conditioned Attention (LCA)}

In the single-law models above, physics enters the transformer in exactly one
place: the residual term of the loss. The law that governs a problem is
\emph{implicit}---the model must infer it from the parameters---and each
domain gets its own network. LCA removes both restrictions: the governing
equation itself becomes part of the input, and one network serves many laws.

\paragraph{Operator vocabulary.} We fix a symbolic vocabulary of
$V=22$ differential and algebraic operators---\texttt{first\_time\_deriv},
\texttt{second\_space\_deriv}, \texttt{self\_advection},
\texttt{moment\_curvature}, \texttt{energy\_conservation},
\texttt{exponential\_relaxation}, and so on (full list in
Supplementary~Table~S1). Each governing equation is expressed as a sparse
binary signature $\mathbf{s}\in\{0,1\}^{V}$ marking the operators it
contains; e.g.\ viscous Burgers $u_t+uu_x=\nu u_{xx}$ gives
\texttt{first\_time\_deriv}$+$\texttt{first\_space\_deriv}$+$
\texttt{second\_space\_deriv}$+$\texttt{self\_advection}$+$
\texttt{diffusion}. The signatures are symbolic facts about the physics,
written by hand from each equation, not learned. Equation~1 of
Supplementary~Table~S1 lists all eight laws with their signatures.

\paragraph{Quantity-labeled input tokens.} The input parameter tokens are
labeled by \emph{physical quantity} from a second shared vocabulary
(\texttt{length}, \texttt{force}, \texttt{modulus}, \texttt{inertia},
\texttt{thickness}, \texttt{speed}, \texttt{angle}, \texttt{stiffness},
\texttt{mass}, \texttt{amplitude}, \texttt{resistance},
\texttt{capacitance}, \texttt{voltage}), so the simply supported beam and
the cantilever---which share the parameter names $L,P,E,I,h$---present the
\emph{identical} token sequence, and the two-parameter domains
(projectile $v_0,\alpha$; pendulum $L,\theta_0$) are distinguished from one
another only by their law signatures. Values are $z$-scored per domain, so
numeric scale cannot leak the domain either. The two information channels
are thus kept strictly separate: \emph{what is given} (quantities) and
\emph{which law governs it} (signature).

\paragraph{Law cross-attention.} An MLP embeds the signature into a law
vector $\mathbf{l}=\mathrm{MLP}(\mathbf{s})\in\mathbb{R}^{d}$; in every
transformer layer the token sequence cross-attends to $\mathbf{l}$
(replicated over the token axis) as keys/values, with a residual and layer
normalization:

\[
\mathbf{h}^{(i+1)} = \mathrm{LN}\big(\mathbf{h}^{(i)} + \mathrm{Attn}(\mathbf{h}^{(i)},\,\mathbf{L},\,\mathbf{L})\big),
\qquad \mathbf{L}=\mathbf{1}\mathbf{l}^{\top},
\]

so attention is conditioned, at every depth, on the differential structure
of the law in force. The same law vector also \emph{gates the shared
readout}: a linear map produces a per-law scale and bias applied to the
final hidden state before the answer and trajectory heads, so the model can
invert per-domain output normalizations ($\log_{10}$ bending answers vs.\
raw periods) directly from the equation. This is a new way to feed physics
into a transformer: not (only) as a loss term on the output, but as an
architectural input that steers both how the network reasons about the
parameters and how it reads them out. For the ablation, the
signature is replaced by a constant vector, so the conditioning stream
exists but carries no equation information; the two conditions are
otherwise identical (same architecture, data, and budget).

\paragraph{Multi-law training.} The multi-law generalist uses one shared
body \emph{and} one shared answer head and trajectory head for all six
50-step domains; each sample's trajectory is normalized per channel by its
own peak (so all heads see comparable $[-1,1]$ targets) and the answer is
$z$-scored per domain, with $\log_{10}$ for the wide-range bending and RC
answers. The physics-consistency residual is still computed per sample from
its \emph{own} governing equation (the sample's domain is known to the
training loop, which is the verifier's job at inference). Training is
otherwise identical to the single-law procedure (AdamW, $\mathrm{lr}=10^{-3}$,
scale-free physics weight ramping in over 10 epochs), with the budget 96
samples per law per seed (576 total) at 250 epochs---37\% of the data a
specialist sees, so the generalist claim is a \emph{data-efficiency} claim
as much as an architecture claim. All LCA results are reported over six
seeds (seed controlling data generation, batch shuffling, and weight
initialization), with paired Wilcoxon tests between the law-conditioned
model and the constant-signature ablation.

\paragraph{Law-swap steering.} To test whether the law vector \emph{causes}
behavior, the signature is swapped at inference on held-out problems whose
parameter tokens are identical across two laws (beam $\leftrightarrow$
cantilever: same five quantity tokens, same sampling ranges). For each sample
$i$ the network is run twice---with the source and target signatures---and
the normalized trajectory errors of the swapped prediction against the two
exact solutions give a steering index
$\mathrm{SI}_i=(d_{i,\mathrm{src}}-d_{i,\mathrm{tgt}})/(d_{i,\mathrm{src}}+d_{i,\mathrm{tgt}})$
(positive: closer to the target law's solution) and a disruption
$d_{i,\mathrm{src}}-d_{i,\mathrm{native}}$ (how far the swap moves the
prediction from the source-law solution). Both are compared between the
law-conditioned model and the constant-signature ablation by paired
Wilcoxon tests over the $n{=}72$ per-condition samples (12 held-out
problems $\times$ 6 seeds); the beam$\to$cantilever direction is reported
because cantilever deflections are the larger family and the normalized
error is free of the scale trap of the reverse direction.

\paragraph{Few-shot adaptation.} The cantilever is held out of pretraining
entirely. A generalist is pretrained on five laws (beam, projectile,
pendulum, spring--mass, RC) at 96 samples per law, 120 epochs, with real or
constant signature; it is then fine-tuned on 24 cantilever samples (25\% of
the 96-sample specialist budget) for 40 epochs at the same learning rate
and physics weight. The from-scratch specialist is trained on the same 24
samples for the same 40 epochs with the same single-law procedure
(256-sample protocol, scaled down). All three conditions are evaluated on
the same 12 held-out cantilever problems (seed $s{+}100$ protocol, never
used in pretraining or fine-tuning); three seeds per condition, with
per-seed values in SI~10.

\paragraph{Physics vs.\ data-only.} On the heat2d domain (2D Poisson field),
identical networks are trained with the physics-consistency loss on
($w_{\mathrm{phys}}=0.05$) or off (data-only, $w_{\mathrm{phys}}=0$), at
64 and 256 samples (25\% and 100\% of the budget), three seeds each, 240
epochs, trajectory weight 3. Evaluation is on each run's own held-out
split (last 1/8 of the generated problems): answer relative error, field
max/mean error of the full 2D prediction against the manufactured exact
solution, and the mean Poisson residual.

\subsection*{Evaluation protocol}

Training/validation split is by generation order (last $\sim$1/8 held out).
Validation metrics: mean relative error
$\frac{1}{N}\sum_i |\hat a_i - a_i|/|a_i|$ on de-scaled answers, and the mean
physics residual on de-scaled trajectories. The held-out generalization
figures (Figures~\ref{fig:predictions}--\ref{fig:burgers}) use 250 freshly
generated problems per domain with seeds distinct from training. The
verification-gate benchmark (Table~\ref{tab:gate}) runs identical missions
with the gate enforced and disabled (\texttt{enforceGate: false} strips
\texttt{expect: ok} so gated failures are not distinguished from
successes), with ground truth computed independently: the benchmark runs
the test suites itself and checks physics missions against the solver's
independent numeric verification. The baseline matrix (Table~\ref{tab:baselines})
trains three kinds (MLP, transformer without physics, transformer with
physics) on each of five domains $\times$ five seeds at identical budgets
(256 samples, 60 epochs, seed controlling data, shuffling, and weight
initialization), and significance is assessed by the two-sided Wilcoxon
signed-rank test on the paired runs---per domain ($n=5$, exact sign-
permutation null) and pooled across domains ($n=25$, asymptotic
$p$)\pn{\citep{wilcoxon1945}}, with Cliff's $\delta$ as the paired effect
size. The external PINN baseline trains \textsc{DeepXDE}\pn{\citep{lu2021}}
(version 1.15.0, PyTorch backend) per instance on three representative
Burgers problems---one fresh FNN per problem (2--50$\times$4--1 tanh),
2{,}000 collocation points, Adam 12k steps then L-BFGS 2k steps (identical
for every instance), exact Cole--Hopf boundary data at
$x=\pm1$---and both models are evaluated with identical metrics on the
model grid. All random seeds, exact hyperparameters, generated datasets,
trained weights, and metrics are committed with the code.

\subsection*{Reproducibility}

The entire system is two packages: a zero-dependency Node.js agent
($\sim$1,600 lines) and a PyTorch physics core ($\sim$3,300 lines), runnable on
a CPU-only laptop with Python $\geq$3.10 and Node $\geq$18. \texttt{npm test}
runs 16 agent tests; \texttt{python3 -m unittest physx.test\_physx} runs 31
physics tests (47 total); both are green in the released commit. The demo
(\texttt{node age.js --demo}) reproduces the end-to-end missions described in
the Results section, and \texttt{npm run bench:gate} reproduces
Table~\ref{tab:gate}. \texttt{python physx/significance.py} reproduces the
statistical tests and \texttt{python physx/deepxde\_baseline.py} (with the
vendored \textsc{DeepXDE} wheel in \texttt{vendor/}) reproduces the
external PINN comparison. \texttt{python physx/run\_matrix.py} reproduces
the full baseline matrix (75 runs).

\section*{Data availability}

All datasets are generated deterministically from the published parameter
ranges and seeds; the generator is included. Trained model weights,
normalization statistics, training logs, journal files, the verification-
gate benchmark results, the significance-test outputs
(\texttt{paper/fig/significance.json}), the \textsc{DeepXDE} comparison
results (\texttt{paper/fig/deepxde\_comparison.json}), and the figure data
(JSON) used to produce Figures 3--11 are included in the repository (DOI
provided on acceptance).

\section*{Code availability}

The complete implementation, tests, and figure-generation scripts are
available at \url{https://github.com/placeholder/age} (DOI
\href{https://doi.org/10.xxxx/xxxxx}{10.xxxx/xxxxx} provided on
acceptance). The code is MIT-licensed.

\section*{Author contributions}

S.R.S. conceived the system, designed and implemented the agent and the
physics core, designed and ran all experiments and benchmarks, and wrote
the manuscript.

\section*{Competing interests}

The authors declare no competing interests. The authors have no affiliation
with Prometheus or any other commercial entity named in this work.

\section*{Acknowledgements}

\textit{Placeholder.} We thank the open-source maintainers of PyTorch, Node.js,
and the physics and agent communities whose tools made this work possible.

\bibliographystyle{unsrtnat}
\begin{thebibliography}{99}
\setlength{\itemsep}{0pt}

\bibitem[LeCun et al.(2015)]{lecun2015} LeCun, Y., Bengio, Y. \& Hinton, G.
Deep learning. \emph{Nature} \textbf{521}, 436--444 (2015).

\bibitem[Silver et al.(2017)]{silver2017} Silver, D. \emph{et al.}
Mastering the game of Go without human knowledge. \emph{Nature}
\textbf{550}, 354--359 (2017).

\bibitem[Mnih et al.(2015)]{mnih2015} Mnih, V. \emph{et al.} Human-level
control through deep reinforcement learning. \emph{Nature} \textbf{518},
529--533 (2015).

\bibitem[Brown et al.(2020)]{brown2020} Brown, T. B. \emph{et al.} Language
models are few-shot learners. In \emph{Advances in Neural Information
Processing Systems 33}, 1877--1901 (2020).

\bibitem[Ouyang et al.(2022)]{ouyang2022} Ouyang, L. \emph{et al.} Training
language models to follow instructions with human feedback. In
\emph{Advances in Neural Information Processing Systems 35} (2022).

\bibitem[Bubeck et al.(2023)]{bubeck2023} Bubeck, S. \emph{et al.} Sparks of
artificial general intelligence: early experiments with GPT-4. Preprint at
\url{https://arxiv.org/abs/2303.12712} (2023).

\bibitem[Wei et al.(2022)]{wei2022} Wei, J. \emph{et al.} Chain-of-thought
prompting elicits reasoning in large language models. In \emph{Advances in
Neural Information Processing Systems 35} (2022).

\bibitem[Vaswani et al.(2017)]{vaswani2017} Vaswani, A. \emph{et al.}
Attention is all you need. In \emph{Advances in Neural Information
Processing Systems 30}, 5998--6008 (2017).

\bibitem[Yao et al.(2023)]{yao2023} Yao, S. \emph{et al.} ReAct: synergizing
reasoning and acting in language models. In \emph{International Conference
on Learning Representations} (2023).

\bibitem[Shinn et al.(2023)]{shinn2023} Shinn, N., Cassano, F., Gopinath, A.,
Narasimhan, K. \& Yao, S. Reflexion: language agents with verbal
reinforcement learning. In \emph{Advances in Neural Information Processing
Systems 36} (2023).

\bibitem[Wang et al.(2023)]{wang2023} Wang, G. \emph{et al.} Voyager: an
open-ended embodied agent with large language models. Preprint at
\url{https://arxiv.org/abs/2305.16291} (2023).

\bibitem[Park et al.(2023)]{park2023} Park, J. S. \emph{et al.} Generative
agents: interactive simulacra of human behavior. In \emph{Proceedings of the
36th Annual ACM Symposium on User Interface Software and Technology} (2023).

\bibitem[Jimenez et al.(2024)]{jimenez2024} Jimenez, C. E. \emph{et al.}
SWE-bench: can language models resolve real-world GitHub issues? In
\emph{International Conference on Learning Representations} (2024).

\bibitem[Li et al.(2022)]{li2022} Li, Y. \emph{et al.} Competition-level code
generation with AlphaCode. \emph{Science} \textbf{378}, 1092--1097 (2022).

\bibitem[Cobbe et al.(2021)]{cobbe2021} Cobbe, K. \emph{et al.} Training
verifiers to solve math word problems. Preprint at
\url{https://arxiv.org/abs/2110.14168} (2021).

\bibitem[Lightman et al.(2024)]{lightman2024} Lightman, H. \emph{et al.}
Let's verify step by step. In \emph{International Conference on Learning
Representations} (2024).

\bibitem[Raissi et al.(2019)]{raissi2019} Raissi, M., Perdikaris, P. \&
Karniadakis, G. E. Physics-informed neural networks: a deep learning
framework for solving forward and inverse problems involving nonlinear
partial differential equations. \emph{J. Comput. Phys.} \textbf{378},
686--707 (2019).

\bibitem[Karniadakis et al.(2021)]{karniadakis2021} Karniadakis, G. E.
\emph{et al.} Physics-informed machine learning. \emph{Nat. Rev. Phys.}
\textbf{3}, 422--440 (2021).

\bibitem[Lu et al.(2021)]{lu2021} Lu, L., Meng, X., Mao, Z. \& Karniadakis,
G. E. DeepXDE: a deep learning library for solving differential equations.
\emph{SIAM Rev.} \textbf{63}, 208--228 (2021).

\bibitem[Chen et al.(2018)]{chen2018} Chen, R. T. Q., Rubanova, Y.,
Bettencourt, J. \& Duvenaud, D. Neural ordinary differential equations. In
\emph{Advances in Neural Information Processing Systems 31} (2018).

\bibitem[Wilcoxon(1945)]{wilcoxon1945} Wilcoxon, F. Individual comparisons
by ranking methods. \emph{Biometrics Bull.} \textbf{1}, 80--83 (1945).

\bibitem[Wang et al.(2021)]{wang2021} Wang, S., Teng, Y. \& Perdikaris, P.
Understanding and mitigating gradient flow pathologies in physics-informed
neural networks. \emph{SIAM J. Sci. Comput.} \textbf{43}, A3055--A3081 (2021).

\bibitem[Takamoto et al.(2022)]{takamoto2022} Takamoto, M. \emph{et al.}
PDEBench: an extensive benchmark for scientific machine learning. In
\emph{Advances in Neural Information Processing Systems 35} (2022).

\bibitem[Kirkpatrick et al.(2017)]{kirkpatrick2017} Kirkpatrick, J.
\emph{et al.} Overcoming catastrophic forgetting in neural networks.
\emph{Proc. Natl Acad. Sci. USA} \textbf{114}, 3521--3526 (2017).

\bibitem[Bommasani et al.(2021)]{bommasani2021} Bommasani, R. \emph{et al.}
On the opportunities and risks of foundation models. Preprint at
\url{https://arxiv.org/abs/2108.07258} (2021).

\bibitem[Sutton(2019)]{sutton2019} Sutton, R. The bitter lesson.
\url{http://www.incompleteideas.net/IncIdeas/BitterLesson.html} (2019).

\bibitem[Temkin(2026)]{techcrunch2026} Temkin, M. Jeff Bezos's Prometheus
raises \$12B to build an `artificial general engineer' for the physical
world. \emph{TechCrunch} (2026).
\url{https://techcrunch.com/2026/06/11/jeff-bezoss-prometheus-raises-12b-to-build-an-artificial-general-engineer-for-the-physical-world/}

\bibitem[Wikipedia contributors(2025)]{wikipedia2025} Wikipedia
contributors. Prometheus (company). \emph{Wikipedia} (2025).
\url{https://en.wikipedia.org/wiki/Prometheus_(company)}

\end{thebibliography}

\end{document}
